%% !TEX program = xelatex
\documentclass[12pt,final]{novelette} % font size, twoside, twocolumn, draft-final, use extarticle for 8-20pt

\title{Diff Eq Midterm 2544}
\author{Kaka Cring}
\date{\today}

\begin{document}

\makeheading{ตัวอย่างข้อสอบ 2301312 สมการเชิงอนุพันธ์ ปรับปรุงจากข้อสอบกลางภาคปลาย ปีการศึกษา 2544}

\subsubsection*{PART A}

\begin{question}
    \item {}
    จงหารอนสเกียนของ $e^x, \cos x$ และ $\sin x$ \\
    และจากรอนสเกียนที่หาได้จงบอกว่าฟังก์ชันทั้งสามเป็นอิสระเชิงเส้นต่อกันหรือไม่บนช่วง $(-\infty, \infty)$ เพราะเหตุใด \qmark{3}

    % ----------------------------------------

    % ----------------------------------------

    \item {}
    จงให้เหตุผลว่าทําไมสมการ $(1-x)y'' + xy' - y = 2x^3 - 6x^2 + 6x$ จึงเป็นปรกติบนช่วง $(1, \infty)$
    และจงหาผลเฉลยบริบูรณ์ของสมการนี้บนช่วง $(1, \infty)$ โดยให้พิจารณาจากฟังก์ชัน $e^x, x$ และ $x^3$ \qmark{8}

    % ----------------------------------------

    % ----------------------------------------

    \item {}
    ลวดสปริงเส้นหนึ่งถูกยึดปลายบนติดกับเพดาน
    เมื่อนําวัตถุหนัก $6$ ปอนด์ มาผูกติดไว้ที่ปลายล่างของลวดสปริงจะทําให้ขดลวดสปริงยืดออก $6$ นิ้ว
    ถ้าดึงวัตถุนี้ลงมาต่ำกว่าตำแหน่งสมดุล $4$ นิ้ว แล้วปล่อยด้วยการผลักวัตถุขึ้นด้วยความเร็ว $\frac83$ ฟุต/วินาที \\
    จงหาสมการของการเคลื่อนที่ของวัตถุพร้อมทั้งบอกค่าแอมพลิจูด และมุมทิ้งช่วง \qmark{10}

    % ----------------------------------------

    % ----------------------------------------

    \item {}
    สมการของการเคลื่อนที่ของวัตถุพร้อมทั้งบอกค่าแอมพลิจูด และมุมทิ้งช่วง \\
    ถ้าเราทราบว่าสมการของการเคลื่อนที่ของวัตถุคือ $x(t) = \frac13 \sin(\pi t - \frac\pi4)$
    \begin{question}
        \item
        จงหา แอมพลิจูด คาบ และความถี่ของการเคลื่อนที่ \qmark{1.5}

        % ----------------------------------------

        % ----------------------------------------

        \item
        จงหาตําแหน่ง ความเร็ว ความเร่งของการเคลื่อนที่ เมื่อเวลา $1$ วินาที หลังจากปล่อยให้วัตถุเกิดการเคลื่อนที่
        พร้อมทั้งอธิบายลักษณะของการเคลื่อนที่ของวัตถุ ณ ขณะนั้น \qmark{3.5}

        % ----------------------------------------

        % ----------------------------------------
    \end{question}
\end{question}

\subsubsection*{PART B}

\begin{question}[resume]
    \item {}
    จงเติมเฉพาะคำตอบลงในช่องว่างที่กำหนดให้โดยไม่ต้องแสดงวิธีทำ \qmark{20}
    \begin{question}
        \item
        ผลเฉลยบริบูรณ์ของสมการ $D(D-3)y = 0$ \qmark{2}

        % ----------------------------------------

        % ----------------------------------------

        \item
        ผลเฉลยบริบูรณ์ของสมการ $(D^2 - 2D + 2)^2 y = 0$ \qmark{2}

        % ----------------------------------------

        % ----------------------------------------

        \item
        ผลเฉลยบริบูรณ์ของสมการสมการเอกพันธ์ุสัมพัทธ์ของสมการ $(D-2)^2(D^2+2D+5)y = e^x$ \qmark{2}

        % ----------------------------------------

        % ----------------------------------------

        \item
        ผลเฉลยเฉพาะหรือปริพันธ์เฉพาะของสมการ $(D-2)^2(D^2+2D+5)y = e^x$ \qmark{2}

        % ----------------------------------------

        % ----------------------------------------

        \item
        ผลเฉลยบริบูรณ์ของสมการสมการเอกพันธ์ุสัมพัทธ์ของสมการ $y^{(4)} + 4y'' = \sin x$ \qmark{2}

        % ----------------------------------------

        % ----------------------------------------

        \item
        ผลเฉลยเฉพาะหรือปริพันธ์เฉพาะของสมการ $y^{(4)} + 4y'' = \sin x$ \qmark{2}

        % ----------------------------------------

        % ----------------------------------------  

        \item
        ตัวดําเนินการเชิงอนุพันธ์อันดับต่ำที่สุดที่ลบล้าง $xe^{2x}$ \qmark{2}

        % ----------------------------------------

        % ----------------------------------------

        \item
        ตัวดําเนินการเชิงอนุพันธ์อันดับต่ำที่สุดที่ลบล้าง $e^{-x}\sin 3x$ \qmark{2}

        % ----------------------------------------

        % ----------------------------------------

        \item
        สมการเชิงอนุพันธ์สามัญเชิงเส้นเอกพันธุ์ที่มีสัมประสิทธิ์เป็นค่าคงตัวอันดับต่ำที่สุดที่มี $e^x$ และ $x^2$ เป็นผลเฉลย \qmark{2}

        % ----------------------------------------

        % ----------------------------------------

        \item
        สมการเชิงอนุพันธ์สามัญเชิงเส้นเอกพันธุ์ที่มีสัมประสิทธิ์เป็นค่าคงตัวอันดับต่ำที่สุดที่มี $x \cos 2x$ และ $x^2 e^x$ เป็นผลเฉลย \qmark{2}

        % ----------------------------------------

        % ----------------------------------------
    \end{question}

    \item {}
    จงหาผลเฉลยเฉพาะของสมการ
    $y''' + y'' = 0; \quad y(0) = 2, \quad y'(0) = 0, \quad y''(0) = 1$ \qmark{5}

    % ----------------------------------------

    % ----------------------------------------
\end{question}

\subsubsection*{PART C}

\begin{question}[resume]
    \item {}
    จงหาปริพันธ์เฉพาะของสมการ $y'' - y = \frac32(1-\cos2x)$ โดยวิธีเทียบสัมประสิทธิ์ \qmark{7}

    % ----------------------------------------

    % ----------------------------------------

    \item {}
    จงใช้ตัวดำเนินการผกผันหาปริพันธ์เฉพาะของสมการ $(D^2 - 9D + 18)y = e^{e^{-3x}}$ \qmark{5}

    % ----------------------------------------

    % ----------------------------------------

    \item {}
    จงใช้ตัวดำเนินการผกผันหาปริพันธ์เฉพาะของสมการ $(D^2+9)y = x \cos x$ \qmark{6}

    % ----------------------------------------

    % ----------------------------------------

    \item {}
    จงหาปริพันธ์เฉพาะของสมการ $(D^3 + D^2 - 2D)y = e^{2x}$ โดยวิธีแปรพารามิเตอร์ \qmark{7}

    % ----------------------------------------

    % ----------------------------------------
\end{question}

\subsubsection*{PART D}

\begin{question}[resume]
    \item {}
    จงหาปริพันธ์เฉพาะของสมการ $(D^2 - 3D + 1)y = x^2 + 15 \sin 2x$ \qmark{8}

    % ----------------------------------------

    % ----------------------------------------

    \item {}
    จงหาผลเฉลยบริบูรณ์ของสมการ $4x^2y'' + 8xy' + y = x^3$ \qmark{6}

    % ----------------------------------------

    % ----------------------------------------

    \item {}
    จงหาผลเฉลยในรูปอนุกรมกำลังรอบจุดสามัญ $x=0$ ของสมการ $y'' - xy' - y = 0$ \\
    (เขียนผลเฉลยให้เห็นอย่างน้อยถึงพจน์ $x^5$) \qmark{10}

    % ----------------------------------------

    % ----------------------------------------
\end{question}

\end{document}